\chapter{Stage 4: Tier Execution}
\label{ch:tiers}

The router has chosen a tier. Stage 4 is where that choice is finally carried out
and an actual answer appears. There are three ways to produce one, and they differ
enormously in cost: read a verified answer straight from memory, generate one from
the model's own knowledge, or retrieve evidence and generate a grounded answer.
This chapter walks all three, the shared reasoning format they obey, and the small
escalation valve that connects the two generating tiers.

\begin{intuition}[Three ways to answer, in rising order of effort]
Faced with a question, a careful person has three moves. If they have answered it
before and remember the answer, they just say it. If they know the area well, they
reason it out from what they already know. And if they are unsure, they look it up
first and then reason from the source. The three tiers are exactly these moves,
made mechanical. The router picked the cheapest move it judged safe; Stage 4
performs it.
\end{intuition}

\section{The shared contract: scaffolded chain of thought}
\label{sec:tiers-contract}

Before the tiers themselves, one convention they share. Both generating tiers, and
the stored episodes that the memory tier serves, follow the same answer format: a
\emph{Reasoning} section followed by a single \emph{Answer} line. This is scaffolded
chain-of-thought~\cite{wei2022chain}, and the uniformity is deliberate and
load-bearing.

\begin{intuition}[Why force the same shape everywhere]
The verifier of \Cref{ch:verifier} reads every answer the same way: it looks for an
explicit chain of reasoning and a cleanly extractable final answer. If a fact-checking
claim, a trivia question, and a multiple-choice item each came back in a different
shape, the verifier would be comparing apples to oranges, and its signals would mean
different things on different benchmarks. Forcing one shape everywhere gives the
verifier a uniform surface to score, which is precisely what lets a single set of
calibrated thresholds govern all five benchmarks. The reasoning scaffold is not just
for the model's benefit; it is the contract the rest of the pipeline reads against.
\end{intuition}

The model is steered into this shape two ways. A system instruction states the format
explicitly and makes the Answer line mandatory, and the generation is \emph{prefilled}
with the literal token ``Reasoning:'' so the model begins inside the scaffold rather
than deciding whether to use it. A representative prompt skeleton, with the
retrieval block that only Tier 3 adds, looks like this.

\begin{lstlisting}[language={}, caption={The scaffolded-CoT prompt skeleton shared by Tiers 2 and 3 (Tier 3 adds the Context block).}, label={lst:prompt}]
<|im_start|>system
Answer with a Reasoning section, then an Answer line. The Answer line is mandatory.
Reasoning: <concise step-by-step thought, 3-4 sentences>
Answer: <final answer in the specified format><|im_end|>
<|im_start|>user
Context:                       # <-- Tier 3 only
[1] <retrieved passage 1>      #     top-5 passages, numbered
...
[5] <retrieved passage 5>
Question: <the query><|im_end|>
<|im_start|>assistant
Reasoning:                     # <-- forced prefix (prefill); the model continues here
\end{lstlisting}

\begin{pitfall}[The mandatory Answer line is a hard-won fix]
An earlier version of the instruction left the Answer line optional, and the cold-start
audit found that on between forty-three and fifty-three percent of samples the model
emitted only a Reasoning section and never reached an Answer line, so the user-facing
result was empty. The silent cost was worse than it first appeared: on the fact-checking
benchmark, about a third of those empty results in fact carried reasoning whose conclusion
already matched the gold verdict, so they were correct answers discarded for a structural
reason rather than a semantic one. The current instruction makes the
Answer line explicitly mandatory and caps the reasoning length so the model reliably
reaches it. The lesson generalises: when a downstream component depends on a
structural feature of the output, that feature has to be demanded in the prompt, not
hoped for. The format is steered further by short per-benchmark worked examples and an
answer-format specification, so a fact-check returns a verdict word, a yes-no question
returns yes or no, and a multiple-choice item returns a single option letter.
\end{pitfall}

\section{Tier 1: serve from memory}
\label{sec:tiers-tier1}

The cheapest tier is barely a computation at all. When the router commits to Tier 1,
the matched episode already holds a verified answer and its trust score, so the tier
reads them straight off the entry and returns. There is no generation and, crucially,
no verification.

\begin{lstlisting}[caption={Tier 1, in full, from \texttt{caem/pipeline.py}.}, label={lst:tier1}]
def _tier1(self, search_with_ids):
    entry, entry_id, similarity = search_with_ids[0]
    return entry.answer, entry.u_stored      # stored answer + its trust score; no model call
\end{lstlisting}

\begin{precise}[Why Tier 1 may skip the verifier]
Skipping verification looks dangerous until we recall what a Tier 1 hit means. The
episode was admitted to memory only because it passed the four-outcome gate of
\Cref{ch:composite} on some earlier cycle, so its answer was \emph{already judged
trustworthy}, and its trust score was written into the entry at that time. Re-running
the nine-signal verifier now would re-derive a verdict the system already holds. So
Tier 1 reads the stored verdict instead of recomputing it, returning in the time of a
single index lookup with no generation, which makes it by construction far cheaper than
the generating tiers; the production design's qualitative target is sub-second serving
rather than a measured per-tier latency, since the batched evaluation logs cannot be
decomposed into a per-tier figure. The cost of judging a fact is paid once, on
the cycle it entered memory, not on every later query that matches it.
\end{precise}

A Tier 1 serve does leave one trace: it updates the episode's usage statistics, the
retrieval count and running success rate of \Cref{sec:memory-feedback}, so that a
heavily served episode's standing can drift with experience and so that a popular
episode is flagged for re-checking. Reading the answer is free; recording that it was
read is the only bookkeeping.

\section{Tier 2: zero-shot generation}
\label{sec:tiers-tier2}

When the router judges the query familiar enough to generate but has no servable
stored answer, it picks Tier 2. The model generates an answer from its own parameters,
conditioned only on the query, with no retrieved evidence. This is where the
self-improvement loop's gains show up at inference time: a query that needed retrieval
in an early cycle can be answered zero-shot once the loop has folded that knowledge into
the adapter. Across the deployed trajectory the zero-shot tier answered its own assigned
queries at least as accurately as the more expensive retrieval tier answered its assigned
queries, with its exact match holding near or above that of the retrieval tier in every
cycle (about $0.57$ by the third cycle against about $0.48$ for the retrieval tier at the
cold start). This is a per-tier-on-assigned-queries comparison, not a claim that retrieval
is globally inferior, and it is precisely the inference-time signature one expects when the
loop has folded retrieval-found knowledge into the model's own parameters.

\begin{lstlisting}[caption={Tier 2 generation, condensed from \texttt{caem/pipeline.py}.}, label={lst:tier2}]
prompt_text, forced_prefix = build_tier2_prompt(query, tokenizer)   # scaffolded ChatML
full_input = prompt_text + forced_prefix                            # prefill "Reasoning:"

output_ids = model.generate(input_ids, max_new_tokens=512,          # cot_max_new_tokens
                            do_sample=False)                        # greedy, reproducible
continuation = tokenizer.decode(output_ids[0, input_len:])          # only the new tokens
answer_str = f"{forced_prefix}{continuation}"                       # re-attach "Reasoning:"
\end{lstlisting}

Three details are worth drawing out. The generation is \emph{greedy}, so the same query
under the same weights yields the same answer, which keeps the whole pipeline
reproducible. Only the newly generated continuation is decoded, and the forced prefix is
re-attached, so the returned string always begins with ``Reasoning:'' and satisfies the
scaffold contract. And the generation budget is capped at five hundred and twelve new
tokens, enough headroom for a few-shot example, a short reasoning section, and the Answer
line. The result is then handed to the verifier of \Cref{ch:verifier}; unlike Tier 1, a
fresh Tier 2 answer has never been judged, so it must be.

\section{Tier 3: retrieval-augmented generation}
\label{sec:tiers-tier3}

The most expensive tier retrieves evidence before generating. It is the architectural
floor, the path the router falls back to whenever the cheaper tiers cannot earn the
query, and the path the classifier of \Cref{ch:ruc} reserves for questions where
grounding genuinely helps.

\begin{fromcode}[The passage store]
Tier 3 retrieves from a read-only store of roughly twenty-one million Wikipedia
passages, each a hundred-word chunk in the dense-passage-retrieval
style~\cite{karpukhin-etal-2020-dense}. The passages are embedded with the \emph{same}
sentence encoder used for episodic memory (\Cref{sec:prereq-embeddings}), so queries,
stored episodes, and corpus passages all live in one $768$-dimensional space. The index
is the approximate inverted-file product-quantised structure of \Cref{sec:prereq-faiss},
sized for the corpus, and it is memory-mapped from disk so that the sixty-four-gigabyte
body stays on disk and only the active cells are paged in. Retrieval uses the high-recall
probe setting, scanning four times as many index cells as the cheap memory-hit
confirmation lookups elsewhere in the system (sixty-four probed cells against sixteen),
because on a Tier 3 query recall matters more than the last millisecond of latency.
\end{fromcode}

The generation itself mirrors Tier 2, with one addition: the top five retrieved passages
are formatted into a numbered Context block above the question, as in the skeleton of
\Cref{lst:prompt}. The retrieve-then-read structure is the standard
retrieval-augmented-generation recipe~\cite{NEURIPS2020_6b493230}, and the model is asked
to reason from the supplied passages rather than from memory alone.

\begin{lstlisting}[caption={Tier 3 retrieve-then-generate, condensed from \texttt{caem/retrieval/rag.py}.}, label={lst:tier3}]
passages = self._retrieve(query, k=5)                       # top-5 from the 21M-passage index
prompt_text, forced_prefix = build_tier3_prompt(query, passages, tokenizer)  # adds Context block
output_ids = model.generate(input_ids, max_new_tokens=512, do_sample=False)  # greedy, grounded
answer = f"{forced_prefix}{tokenizer.decode(output_ids[0, input_len:])}"      # scaffold contract
\end{lstlisting}

Like a Tier 2 answer, a Tier 3 answer is fresh and unjudged, so it passes through the
verifier and the gate, and a Tier 3 answer that verifies well is stored as a new episode,
so that a future similar query can be served from the cheaper tiers. This is the
mechanism by which the compute ladder \emph{lowers} over time: yesterday's expensive
retrieval becomes tomorrow's memory hit.

\section{The escalation edge}
\label{sec:tiers-escalation}

One connection runs directly between the two generating tiers. If a Tier 2 generation
comes back empty, or the generation call fails outright, the pipeline does not return an
empty answer; it quietly escalates that single query to Tier 3 and marks the result as
escalated.

\begin{pitfall}[The reported tier still reads 2 after an escalation]
The escalation is a robustness valve, not a re-routing. When it fires, the answer comes
from Tier 3, but the routing decision is not rewritten: the recorded tier stays at two,
with a separate flag noting that the query escalated. Anyone reading the logs should know
that an escalated Tier 2 result was in fact produced by retrieval. The valve exists
because an empty zero-shot answer is worse than a grounded one, so a failed cheap attempt
falls forward to the expensive path rather than failing the query.
\end{pitfall}

\section{What leaves Stage 4, and where it goes}
\label{sec:tiers-exit}

\Cref{tab:tiers} collects the three tiers side by side, and the single most important
column is the last one. A Tier 1 answer carries a trust score already and goes straight to
the output stage. A Tier 2 or Tier 3 answer is a fresh, unjudged string, and it goes to the
verifier of \Cref{ch:verifier} before the system will trust it, store it, or refuse it.

\begin{table}[htbp]
\centering
\caption[The three tiers compared]{The three execution tiers. Only the two generating
tiers produce a fresh answer that must be verified; the memory tier serves an answer that
was already judged on a past cycle.}
\label{tab:tiers}
\small
\renewcommand{\arraystretch}{1.25}
\begin{tabularx}{\textwidth}{@{}l >{\raggedright\arraybackslash}X c c@{}}
\toprule
\textbf{Tier} & \textbf{Mechanism} & \textbf{Generates?} & \textbf{Verified now?} \\
\midrule
1 (memory-direct) & read the stored answer and its trust score off the matched episode
  & no & no (judged on a past cycle) \\
2 (zero-shot) & generate from the model's own knowledge, query only
  & yes & yes \\
3 (retrieval) & retrieve top-5 passages from the 21M-passage corpus, then generate grounded
  & yes & yes \\
\bottomrule
\end{tabularx}
\end{table}

\begin{bythenumbers}[Tier-execution constants]
\begin{itemize}[leftmargin=1.3em, itemsep=1pt]
  \item Generation: greedy (deterministic) on the frozen backbone plus its adapter; cap $512$ new tokens.
  \item Forced prefix: every generated answer is prefilled with ``Reasoning:'' to enter the scaffold.
  \item Tier 3 retrieval: top $5$ passages from $\approx 21$ million Wikipedia chunks; one shared $768$-dimensional embedding space; high-recall probe setting; index memory-mapped from disk.
  \item Escalation: an empty or failed Tier 2 generation falls forward to Tier 3, flagged as escalated.
  \item Verification: Tiers 2 and 3 go to the verifier; Tier 1 does not.
\end{itemize}
\end{bythenumbers}

\begin{defence}[If an examiner asks: ``why not verify Tier 1 answers too, to be safe?'']
Two reasons, one about correctness and one about cost. On correctness, a Tier 1 answer was
admitted to memory only by passing the gate, and the memory is re-judged in full at every
\emph{committed} cycle boundary by the retroactive re-verification of \Cref{ch:cycle}, so a
stored answer is rarely more than a cycle old without a fresh verdict. On cost, the entire value of the memory
tier is that it answers recurring queries for the price of a lookup; verifying every Tier 1
hit would erase that saving and make the cheapest tier nearly as expensive as the others.
The design verifies a fact when it enters memory and re-verifies it each cycle, rather than
on every serve, which is where the asymmetry of the architecture pays off.
\end{defence}

Every tier has now produced an answer, and for the two generating tiers that answer is still
unjudged. \Cref{ch:verifier} opens Stage 5, the nine-signal verifier, the component that
decides how much to trust a freshly generated answer.
